\documentclass[11pt, a4paper]{article}

% --- UNIVERSAL PREAMBLE BLOCK ---
\usepackage[a4paper, top=2.5cm, bottom=2.5cm, left=2cm, right=2cm]{geometry}
\usepackage{fontspec}
\usepackage[spanish, bidi=basic, provide=*]{babel}

\babelprovide[import, onchar=ids fonts]{spanish}
\babelprovide[import, onchar=ids fonts]{english}

% Configuración de la fuente Noto Sans
\babelfont{rm}{Noto Sans}

% Paquetes técnicos
\usepackage{amsmath}
\usepackage{booktabs}
\usepackage{graphicx}
\usepackage{listings}
\usepackage{xcolor}
\usepackage[siunitx, american]{circuitikz}
\usepackage{tikz}
\usetikzlibrary{shapes.geometric, arrows, positioning, calc, perspective}
\usepackage{enumitem}
\usepackage{hyperref}
\usepackage{array}
\usepackage{caption}

% --- ESTILO DE CÓDIGO ---
\lstset{
    language=Python,
    basicstyle=\ttfamily\tiny,
    keywordstyle=\color{blue},
    stringstyle=\color{red},
    commentstyle=\color{green!60!black},
    breaklines=true,
    numbers=left,
    numberstyle=\tiny\color{gray},
    frame=single,
    backgroundcolor=\color{gray!5}
}

% --- ESTILOS TIKZ ---
\tikzset{
    startstop/.style={rectangle, rounded corners, minimum width=3cm, minimum height=1cm, text centered, draw=black, fill=red!20, font=\small},
    process/.style={rectangle, minimum width=4.2cm, minimum height=1.5cm, text width=4cm, text centered, draw=black, fill=orange!15, font=\small},
    decision/.style={diamond, aspect=2.2, minimum width=3.8cm, minimum height=1cm, text centered, draw=black, fill=green!15, font=\small},
    arrow/.style={thick,->,>=stealth}
}

\begin{document}

% --- PORTADA ---
\begin{titlepage}
    \centering
    {\LARGE Universidad Autónoma de Nuevo León}\\[0.5cm]
    {\Large Facultad de Ingeniería Mecánica y Eléctrica}\\[2cm]
    
    {\Large Laboratorio de Técnicas de Medida en Aeronáutica}\\[1cm]
    
    {\Huge Práctica 6}\\[0.5cm]
    {\Large Velocimetría de Partículas por Imagen (PIV) y Flujo Óptico}\\[2cm]
    
    Elaborado para el ciclo escolar:\\
    2026\\[2cm]
    
    \vfill
    
    Resumen Ejecutivo\\
    Esta sesión aborda la técnica PIV para la obtención de campos de velocidad bidimensionales instantáneos. A diferencia de las mediciones de punto único, el estudiante analizará pares de imágenes sintéticas utilizando algoritmos de flujo óptico (Farnebäck) para reconstruir la dinámica global del fluido. El enfoque analítico se centra en la identificación de estructuras coherentes, como vórtices de Rankine, y el cálculo de la vorticidad local para diagnosticar el comportamiento del flujo en componentes aeroespaciales.
\end{titlepage}

\tableofcontents
\newpage

% --- NOMENCLATURA ---
\section*{Nomenclatura}
\begin{description}
    \item[$\Delta t$] Intervalo de tiempo entre pulsos láser [s]
    \item[$\Delta x, \Delta y$] Desplazamiento de partículas en el plano [px o m]
    \item[$u, v$] Componentes cartesianas del vector de velocidad [m/s]
    \item[$|V|$] Magnitud de la velocidad local [m/s]
    \item[$\omega_z$] Vorticidad (componente normal al plano) [1/s]
    \item[$\psi$] Función de corriente (Streamlines)
    \item[CCD/CMOS] Dispositivos de captura de imagen digital
\end{description}

\section{Introducción}
La Velocimetría por Imágenes de Partículas (PIV) es una técnica óptica no intrusiva fundamental para caracterizar flujos complejos. Su capacidad para proporcionar una visión global instantánea permite estudiar fenómenos como el desprendimiento de capa límite en perfiles alares y la interacción de estelas en motores de turbina. En esta práctica, el estudiante procesará pares de imágenes de partículas trazadoras utilizando herramientas de visión por computadora para derivar campos vectoriales y escalares de interés aerodinámico.

\section{Objetivos de Aprendizaje}
\begin{itemize}[label=-]
    \item Comprender los principios de iluminación por lámina láser y sincronización temporal en PIV.
    \item Implementar algoritmos de flujo óptico denso para la estimación de desplazamientos de partículas.
    \item Visualizar campos vectoriales mediante mapas de flechas (quiver plots) y líneas de corriente.
    \item Calcular magnitudes derivadas como la magnitud de velocidad y la vorticidad local.
    \item Interpretar la física del flujo a partir de la identificación de centros de rotación y zonas de cizalladura.
\end{itemize}

\section{Fundamento Teórico}

\subsection{Principio de Funcionamiento de PIV}
El sistema PIV se basa en la captura de dos estados sucesivos de un flujo sembrado con trazadores. Una lámina láser ilumina un plano específico en dos instantes ($t$ y $t+\Delta t$). El desplazamiento promedio de grupos de partículas en ventanas de interrogación define el vector de velocidad local:
\begin{equation}
    \vec{V} = \lim_{\Delta t \to 0} \frac{\Delta \vec{x}}{\Delta t}
\end{equation}

\begin{figure}[htbp]
    \centering
    \begin{tikzpicture}[scale=0.8, x={(0.866cm,0.5cm)}, y={(-0.866cm,0.5cm)}, z={(0cm,1cm)}]
        % Túnel de viento (sección)
        \draw[thick] (0,0,0) -- (4,0,0) -- (4,4,0) -- (0,4,0) -- cycle;
        \draw[thick] (0,0,2) -- (4,0,2) -- (4,4,2) -- (0,4,2) -- cycle;
        \draw[thick] (0,0,0) -- (0,0,2);
        \draw[thick] (4,0,0) -- (4,0,2);
        \draw[thick] (4,4,0) -- (4,4,2);
        \draw[thick] (0,4,0) -- (0,4,2);

        % Lámina Láser (Sheet)
        \fill[green, opacity=0.3] (0,2,0) -- (4,2,0) -- (4,2,2) -- (0,2,2) -- cycle;
        \draw[green!70!black, thick] (0,2,0) -- (4,2,2);
        \node at (4.5,2,2.5) [font=\tiny, color=green!50!black] {Lámina Láser};

        % Cámara
        \draw[fill=gray!30] (2,5,1) rectangle (2.5,5.5,1.5);
        \draw[thick] (2.25,5,1.25) -- (2.25,3,1); % Eje óptico
        \node at (2.25,6,1.25) [font=\tiny] {Cámara CCD};

        % Flujo
        \draw[blue, ->, thick] (-1,2,1) -- (0.5,2,1);
        \node at (-1.5,2,1) [font=\tiny] {Flujo con Partículas};
    \end{tikzpicture}
    \caption{Configuración experimental típica para PIV en un plano de flujo.}
\end{figure}

\subsection{Vorticidad y Análisis Escalar}
La vorticidad ($\omega_z$) cuantifica la rotación local del fluido. En un campo bidimensional, se obtiene a partir del gradiente de las componentes de velocidad:
\begin{equation}
    \omega_z = \frac{\partial v}{\partial x} - \frac{\partial u}{\partial y}
\end{equation}
Valores positivos indican una rotación antihoraria, mientras que valores negativos corresponden a una rotación horaria, permitiendo localizar el núcleo de vórtices o estructuras turbulentas.

\section{Metodología de Análisis de Datos}

\subsection{Flujo de Procesamiento Algorítmico}
El análisis se realiza mediante un script de Python que utiliza la biblioteca OpenCV para implementar el flujo óptico de Farnebäck.

\begin{figure}[htbp]
    \centering
    \begin{tikzpicture}[node distance=1.8cm, auto]
        \node (load) [startstop] {Carga de Imágenes ($I_1, I_2$) en Grises};
        \node (flow) [process, below=of load] {Cálculo de Flujo Óptico (Farnebäck)};
        \node (uv) [process, below=of flow] {Separación de Componentes $u(x,y)$ y $v(x,y)$};
        \node (vort) [process, below=of uv] {Cálculo de Vorticidad vía \texttt{np.gradient}};
        \node (viz) [process, below=of vort] {Visualización: Quiver, Magnitud y Vorticidad};
        \node (rep) [startstop, right=of viz, xshift=2cm] {Informe AIAA};
        
        \draw [arrow] (load) -- (flow);
        \draw [arrow] (flow) -- (uv);
        \draw [arrow] (uv) -- (vort);
        \draw [arrow] (vort) -- (viz);
        \draw [arrow] (viz) -- (rep);
    \end{tikzpicture}
    \caption{Secuencia de procesamiento para la estimación del campo de velocidad.}
\end{figure}

\section{Casos de Estudio y Datasets}

\subsection{Descripción de las Imágenes Sintéticas}
Se dispone de pares de imágenes que representan un Vórtice de Rankine en diferentes instantes. El estudiante debe seleccionar un par coherente para el análisis.

\begin{table}[htbp]
    \centering
    \begin{tabular}{lllc}
        \toprule
        ID Experimento & Imagen $t_0$ & Imagen $t_0 + \Delta t$ & Tipo de Flujo \\
        \midrule
        Caso A & rankine\_vortex03\_0.jpg & rankine\_vortex03\_1.jpg & Vórtice Simple \\
        Caso B & rankine\_vortex02\_1.jpg & rankine\_vortex03\_0.jpg & Evolución Temporal \\
        \bottomrule
    \end{tabular}
    \caption{Datasets disponibles para el análisis PIV mediante flujo óptico.}
\end{table}

\section{Actividades a Realizar}
\begin{enumerate}
    \item Procesamiento de Imágenes: Cargue el par seleccionado en escala de grises y verifique que las dimensiones coincidan.
    \item Estimación de Campo Denso: Ejecute el algoritmo de Farnebäck ajustando los niveles de la pirámide y el tamaño de la ventana (\texttt{winsize}) para minimizar el ruido.
    \item Generación de Mapas Vectoriales: Cree un gráfico de tipo \textit{quiver} submuestreando la imagen (ej. un vector cada 20 píxeles) para facilitar la interpretación visual.
    \item Análisis de Vorticidad: Identifique el centro del vórtice localizando el punto de máxima magnitud de vorticidad y discuta su signo.
\end{enumerate}

\section{Análisis de Resultados y Graficación}
En el reporte AIAA se deben analizar las siguientes situaciones:
\begin{enumerate}
    \item Campo de Magnitud: Grafique la rapidez $|V|$ como un mapa de calor. Identifique si existe un núcleo de rotación rígida (velocidad proporcional al radio) y una zona de decaimiento externa.
    \item Visualización de Vorticidad: Utilice un mapa de colores divergente (blue-white-red). Explique la presencia de picos de vorticidad y su relación con el centro de la estructura.
    \item Sensibilidad del Algoritmo: Discuta cómo varía el campo vectorial al cambiar el parámetro \texttt{poly\_sigma} en el script. ¿Cómo afecta el suavizado a la detección de pequeños vórtices?
\end{enumerate}

\section{Preguntas de Discusión Electrónica y Física}
\begin{enumerate}
    \item ¿Qué representa físicamente el intervalo $\Delta t$ en el contexto de la escala de velocidades de la aeronave y la resolución de la cámara?
    \item Explique por qué el uso de imágenes sintéticas limpias facilita el análisis de flujo óptico frente a imágenes reales de túnel de viento con reflexiones en las paredes.
    \item En la instrumentación PIV, ¿cuál es el propósito de sincronizar el disparo del láser con la apertura del obturador de la cámara mediante un controlador de retardos?
    \item ¿Cómo influye el tamaño de las partículas trazadoras en su capacidad para seguir las fluctuaciones de alta frecuencia de la turbulencia?
\end{enumerate}

\section{Lineamientos del Reporte (AIAA)}
El reporte debe presentar los tres mapas generados (Vectores, Magnitud y Vorticidad) con escalas de ejes en píxeles. Se debe incluir una interpretación física del vórtice analizado, comparándolo con el modelo teórico de un vórtice de Rankine.

\newpage
\appendix
\section{Apéndice: Script de Procesamiento PIV (\texttt{analisis\_piv.py})}
Este script automatiza el cálculo del flujo óptico y la generación de magnitudes derivadas.

\begin{lstlisting}[language=Python]
import cv2
import numpy as np
import matplotlib.pyplot as plt

def procesar_piv(path1, path2):
    # Carga de imagenes
    img1 = cv2.imread(path1, cv2.IMREAD_GRAYSCALE)
    img2 = cv2.imread(path2, cv2.IMREAD_GRAYSCALE)
    
    # Calculo de Flujo Optico (Farneback)
    flow = cv2.calcOpticalFlowFarneback(img1, img2, None, 0.5, 3, 15, 3, 5, 1.2, 0)
    u = flow[..., 0]
    v = flow[..., 1]
    
    # Magnitud y Vorticidad
    mag = np.sqrt(u**2 + v**2)
    dy, dx = np.gradient(v)
    duy, dux = np.gradient(u)
    vort = dx - duy # dwz = dv/dx - du/dy
    
    return u, v, mag, vort

# Ejemplo de visualizacion Quiver
# plt.quiver(x[::20,::20], y[::20,::20], u[::20,::20], v[::20,::20])
\end{lstlisting}

\section{Referencias}
\begin{enumerate}
    \item Raffel, M. et al. (2018). Particle Image Velocimetry: A Practical Guide. Springer.
    \item Adrian, R. J. (1991). Particle-imaging techniques for experimental fluid mechanics.
    \item Documentación oficial OpenCV - Optical Flow.
\end{enumerate}

\end{document}